\section{Créditos}
Este documento es una recopilación hecha por la comunidad. Gracias a todas las personas de la UNAM que cooperaron aportando y compartiendo sus códigos: cada carrera cuenta. Documento original mantenido por \textbf{Etmis} y colaboradores, en esta versión de latex también puede contribuir cualquier persona interesada y si no les gusta a los creadores originales lo borro.

\begin{myquote}{Enlaces originales}
Hoja de códigos (Google Docs):\\
\url{https://docs.google.com/document/d/1o8UgQYTTjpBznUktuLK0WDKmX9nV4Ao2UIWs1MJQ2-s/edit?usp=sharing}

\vspace{0.4em}
Hoja de cálculo (Google Sheets):\\
\url{https://docs.google.com/spreadsheets/d/1SvWsZoNucry4_kdDRg472s3KdWpfSO0qfrXBOsQd2rA/edit?usp=sharing}
\end{myquote}

\section{Cómo contribuir}
\begin{itemize}[leftmargin=*]
  \item No borres los códigos que ya estén publicados. Ayuda a conservar el trabajo de la comunidad.
  \item Puedes agregar tu código aunque ya exista el de tu carrera o plantel.
  \item Separa los códigos nuevos con una coma.
  \item Este documento está separado por entidades: CCH, preparatorias, ENES, FES, facultades de CU y otras entidades.
  \item Verifica que el código corresponda a tu carrera: algunos códigos no funcionan para la carrera que dicen. Puedes reportar errores en el repositorio de GitHub, abre un \emph{issue} o propón el cambio en el archivo \texttt{reportes.tex}. Si no eres colaborador, GitHub creará un \emph{fork} y un \emph{pull request}. Esto se explica mejor en el README.
  \item Si un plantel o carrera no tiene código todavía, aparece sin códigos en su sección.
\end{itemize}

\section{Acerca de este documento}
Este repositorio conserva los códigos en formato \LaTeX{}, con un archivo por plantel, de modo que cualquier persona pueda proponer cambios sin que se pierda la información. También se incluyen los apartados especiales de \emph{Tocayo}, \emph{Casi Familia}, \emph{Toda la república}, \emph{Medallas Secretas}, \emph{Pokemones}, \emph{Pastel Compartido} y \emph{Compa de Banca}.
